% Created 2026-08-29 za 20:20
% Intended LaTeX compiler: pdflatex
\documentclass[11pt]{article}
\usepackage[utf8]{inputenc}
\usepackage[T1]{fontenc}
\usepackage{graphicx}
\usepackage{longtable}
\usepackage{wrapfig}
\usepackage{rotating}
\usepackage[normalem]{ulem}
\usepackage{amsmath}
\usepackage{amssymb}
\usepackage{capt-of}
\usepackage{hyperref}
\date{\today}
\title{}
\hypersetup{
 pdfauthor={},
 pdftitle={},
 pdfkeywords={},
 pdfsubject={},
 pdfcreator={Emacs 30.2 (Org mode 9.7.11)}, 
 pdflang={English}}
\begin{document}

\newif\ifisdraft\isdraftfalse
\newif\ifprint\printfalse
\%\documentclass[twoside]{article}

\% Plots
\usepackage[x11names,svgnames,table]{xcolor}
\usepackage{tikz,pgfplots,pgfplotstable}
\pgfplotsset{width=0.75\linewidth, compat=1.18}
\usepgfplotslibrary{external}
\tikzexternalize[prefix=tikz/, optimize=true, optimize command away=\includepdf]

\tikzifexternalizing\{\%
    \% don't include package XYZ here
\}\{\%
    \usepackage{pdfpages}
\}\%

\% Glossaries
\usepackage[acronym]{glossaries}

\makenoidxglossaries
\loadglsentries{~/Documents/notes/global/begrippenlijst.tex}
\%\renewcommand*{\seename}{zie}

\usepackage{lineno}



\% Quotes work now
\usepackage[autostyle]{csquotes}

\% Footers and page numbers  
\usepackage{fancyhdr}

\% Changepage allows stuff that doesn't fit to go off the page with the \begin{adjustwidth}{-2cm}{} command. it also helps with the intentionally blank page on odd page numbers.
\usepackage[strict]{changepage}

\% Colour section headings
\usepackage{xcolor}

\% Stop sectsty from making warnings that don't make sense
\usepackage[immediate]{silence}
\WarningFilter[temp]{latex}{Command} \% silence the warning
\usepackage{sectsty}
\DeactivateWarningFilters[temp] \% So nothing unrelated gets silenced

\makeatletter \% disable the runtime redefinitions
\let\SS@makeulinesect\relax
\let\SS@makeulinepartchap\relax
\makeatother

\% More symbols
\usepackage{textcomp}

\% Ifthenelse
\usepackage{ifthen}

\% Needed for ununderstood reasons
\usepackage{import}

\% Part of adding red text to todo items list
\usepackage{nameref}

\% --- Manual To-Do List Setup ---
\makeatletter

\% Setup infrastructure to disable errors when a list is empty since it is sometimes in some runs of lualatex
\let$\backslash$@oldnoitemerr$\backslash$@noitemerr \%Save the command definition                      
\newcommand\noitemerroroff{\let\@noitemerr\relax}
\newcommand\noitemerroron{\let\@noitemerr\@oldnoitemerr}

\% 1. Create a new file handle for our to-do list
\newwrite\todofile

\% 2. Open the .tdo file when the document begins.
\%    \jobname refers to the main name of your .tex file.
\%\AtBeginDocument{\openout\todofile=\jobname.tdo}
\immediate\openout\todofile=output.tdo

\newcounter{todocounter}

\% 4. Define the command that will format each item in the list.
\%    This command will be written to the .tdo file and executed later.
\%    \#1 is the to-do text, \#2 is the page number.
\%    \dotfill adds dots before the page number for alignment.
\newcommand{\todolistitem}[3]\{\hyperlink{#3}\{\item \#1 \dotfill \textit{p.~#2}\}\}

\% 5. Define your \todo command.
\newcommand{\todo}[1]\{\%
  \stepcounter{todocounter}\%
  \hypertarget{\thetodocounter}{}\%
  \protected@write\todofile{}\{\%
    \% Add $\backslash$@currentlabelname for debugging
    \protect\todolistitem{#1}{\thepage}\{\the\value{todocounter}\}\%
  \}\%
\}

\% 6. Define the \listoftodos command to display the list.
\newcommand{\listoftodos}\{\%
  \section*{List of To-Do Items}\% Set a title
  \noitemerroroff
  \begin{itemize}\% Use an itemized list environment
  \% Check if the .tdo file exists. If yes, input it.
  \% If not (e.g., first run), print a message.
  \IfFileExists{output.tdo}\{\%
      \input{|"sh ~/Documents/notes/global/sortscript.sh"}
      \% $\backslash$@input\{\jobname.tdo\}\% Read the .tdo file
  \}\{\%
      \item No To-Do items yet!\%
  \}\%

  \end{itemize}\% Close the list environment
  \noitemerroron
\}

\% Part of adding red text to todo items list
\newcommand*{\currentname}{\@currentlabelname}

\makeatother

\usepackage{luacolor,lua-ul} \%for usage of style attributes - background color

\ifprint
\% \newcommand{\fout}[2][_VAR_ Rode tekst]\{\ifthenelse\{=\{\#1\}\{notodo\}\}\{\}\{\todo{\theparagraph: #1}\}\textit{#2}\}
\% \newcommand{\highlight}[1]\{\todo{\theparagraph: #1}\colorbox{lightgray}{#1}\}
\newcommand{\fout}[2][_VAR_ Rode tekst]{}
\newcommand{\highlight}[1]{}
\else
\newcommand{\fout}[2][_VAR_ Rood (\immediate\currentname)]\{\ifthenelse\{=\{\#1\}\{notodo\}\}\{\}\{\todo{\theparagraph: #1}\}\textcolor{red}{#2}\}
\newcommand{\highlight}[1]\{\todo{\theparagraph: #1}\highLight{#1}\}
\fi

\% Watermark
\% \usepackage{draftwatermark}
\% \SetWatermarkText{TECHNASIUM}
\% \SetWatermarkScale{0.6}
\% \SetWatermarkLightness{0.9}

\% Frontmatter, Mainmatter, and Backmatter
\makeatletter

\newcommand\frontmatter\{\%
  \%$\backslash$@mainmatterfalse
  \pagenumbering{roman}
  \}

\newcommand\mainmatter\{\%
 \% $\backslash$@mainmattertrue
  \pagenumbering{arabic}
  \}

\newcommand\backmatter\{\%
  \if@openright
    \cleardoublepage
  \else
    \clearpage
  \fi
 \% $\backslash$@mainmatterfalse
   \}

\makeatother

\% language
\usepackage{polyglossia}
\%\setdefaultlanguage[babelshorthands,tremahyphenation]{dutch}
\setdefaultlanguage[variant=british,ordinalmonthday]{english}

\% Page geometry
\usepackage[a4paper,margin=2cm, marginparwidth=1.6cm,marginparsep=0.2cm]{geometry}

\% Verbatim
\usepackage{verbatim}

\% Math
\usepackage{amsmath}

\% Make H float specifier which makes images work in multicol environments
\usepackage{float}

\% Alter list spacing
\usepackage{enumitem}
\setlist{nosep} \% or \setlist{noitemsep} to leave space around whole list

\% Symbolen. Vooral graden
\usepackage{gensymb}

\% Images
\usepackage{graphicx}
\% Typography
\usepackage{microtype}
  \% Quotes
\usepackage{csquotes}
  \% Fonts
  \usepackage[T1]{fontenc}
  \usepackage{wedn}
  \%Or use Times?
  \%\%\usepackage{mathptmx}
  \% Baskervaldx
  \iftrue
   \usepackage[lf]{Baskervaldx} \% lining figures
   \usepackage[bigdelims,vvarbb]{newtxmath} \% math italic letters from Nimbus Roman
   \usepackage[cal=boondoxo]{mathalfa} \% mathcal from STIX, unslanted a bit
   \renewcommand*\oldstylenums[1]\{\textosf{#1}\}
  \fi





\% Bibliography
\usepackage{[}
backend=biber,
style=apa,
]\{biblatex\}
\addbibresource{~/Documents/notes/global/bronnen.bib}
\DeclareLanguageMapping{latin}{english}

\% \#E68A1E, \#F0AC1A, \#F6C521, \#FBE106

\% h1 kleur
\definecolor{Accent1}{HTML}{2E3440}

\% h2 kleur
\definecolor{Accent2}{HTML}{3B4252}

\% h3 kleur
\definecolor{Accent3}{HTML}{434C5E}

\% h4 kleur
\definecolor{Accent4}{HTML}{4C566A}

\sectionfont\{\color{Accent1}\}  \% sets colour of sections
\subsectionfont\{\color{Accent2}\}
\subsubsectionfont\{\color{Accent3}\}

\% Don't ask how this works, I don't really know either. In theory it makes paragraphs behave as subsubsubsections though and in practice too!
\setcounter{secnumdepth}{4}
\setcounter{tocdepth}{4}

\makeatletter
\renewcommand{\paragraph}\{
$\backslash$@startsection\{paragraph\}\{4\}\{\z@\}\%
                                    \{3.25ex $\backslash$@plus1ex $\backslash$@minus.2ex\}\%
                                    \{0.1ex\}\%
                                    \{\normalfont\small\bfseries\color{Accent4}\}
\}

\% Alias subsubsubsection to paragraph
\newcommand{\subsubsubsection}{\paragraph}

\makeatother

\% For writing down notes and stuff
\usepackage{[}
linenumber,
incolumn,
]\{mindflow\}

\% Appendix
\usepackage[title,titletoc]{appendix}

\% For code formatting
\usepackage{listings}

\% Better references
\ifprint
\usepackage[colorlinks=false, allcolors=blue]{hyperref}
\else
\usepackage[colorlinks=true, allcolors=blue]{hyperref}
\fi
\usepackage[nameinlink]{cleveref}



\% Misc
\let\tqs\textquotesingle

\% Multiple columns
\usepackage{multicol}

\% We want stuff like the TOC and the bibliography included in the TOC
\usepackage{tocbibind}

\% Gifs
\usepackage{animate}

\% Tables
\usepackage{booktabs}

\% Tabel alignment
\usepackage{array}
\newcolumntype{P}[1]\{>\{\centering\arraybackslash\}p\{\#1\}\}

\% Add a colon at a description list
\renewcommand*{\descriptionlabel}[1]\{\hspace{\labelsep}\textbf{#1:}\}

\usepackage{lipsum}
\setotherlanguage{latin}

\usepackage{titling}

\% lipsum rood maken
\NewCommandCopy{\oldlipsum}{\lipsum}
\renewcommand{\lipsum}[1][1]\{\fout[_VAR_ Lipsum (\immediate\currentname)]{\oldlipsum[#1]}\}

\makeatletter
\def\fulltitle{\subject --- \thetitle --- \thedate}
\makeatother


\newcommand{\NotesDocumentSetup}\{\%
\pagestyle{fancy}
\fancyhead{} \% clear all header fields
\setlength{\headheight}{32.04134pt}
\addtolength{\topmargin}{-15.04134pt}
\fancyhead[CO,CE]\{\textcolor{Accent1}\{\textbf{\fulltitle}\}\}
\%\fancyhead[RO,LE]\{\includegraphics[width=1cm]{Images/SolarTeamEindhovenLogoSquare.png}\}
\%\fancyhead[LO,RE]\{\includegraphics[width=1cm]{Images/GooisLyceumLogoSquare.png}\}
\fancyfoot{} \% clear all footer fields
\fancyfoot[LE,RO]{\thepage}

\%\tikzifexternalizing\{\%
\%\}\{\%
\%    \includepdf{Images/PWSVoorblad.pdf}
\%\}\%

\frontmatter

\setcounter{page}{1}
\maketitle

\clearpage
\clearpage

\clearpage
\tableofcontents
\clearpage

\setcounter{page}{1}
\mainmatter
\}

\AtBeginDocument{\NotesDocumentSetup}
\end{document}
